\documentclass[10pt,a4paper]{amsart}
\usepackage{amsmath,amssymb,amsthm,graphicx}
\usepackage{verbatim}
\newenvironment{reference}{\comment}{\endcomment}
%\newenvironment{reference}{}{}
\newenvironment{slogan}{\comment}{\endcomment}
\newenvironment{history}{\comment}{\endcomment}

% For commutative diagrams we use Xy-pic
\usepackage[all]{xy}

% We use 2cell for 2-commutative diagrams.
\xyoption{2cell}
% Theorem environments.
%
\theoremstyle{plain}
\newtheorem{theorem}[subsection]{Theorem}
\newtheorem{proposition}[subsection]{Proposition}
\newtheorem{lemma}[subsection]{Lemma}

\theoremstyle{definition}
\newtheorem{definition}[subsection]{Definition}
\newtheorem{example}[subsection]{Example}
\newtheorem{exercise}[subsection]{Exercise}
\newtheorem{situation}[subsection]{Situation}

\theoremstyle{remark}
\newtheorem{remark}[subsection]{Remark}
\newtheorem{remarks}[subsection]{Remarks}

\numberwithin{equation}{subsection}

% Macros
%
\def\lim{\mathop{\mathrm{lim}}\nolimits}
\def\colim{\mathop{\mathrm{colim}}\nolimits}
\def\Spec{\mathop{\mathrm{Spec}}}
\def\Hom{\mathop{\mathrm{Hom}}\nolimits}
\def\Ext{\mathop{\mathrm{Ext}}\nolimits}
\def\SheafHom{\mathop{\mathcal{H}\!\mathit{om}}\nolimits}
\def\SheafExt{\mathop{\mathcal{E}\!\mathit{xt}}\nolimits}
\def\Sch{\mathit{Sch}}
\def\Mor{\mathop{\mathrm{Mor}}\nolimits}
\def\Ob{\mathop{\mathrm{Ob}}\nolimits}
\def\Sh{\mathop{\mathit{Sh}}\nolimits}
\def\NL{\mathop{N\!L}\nolimits}
\def\CH{\mathop{\mathrm{CH}}\nolimits}
\def\proetale{{pro\text{-}\acute{e}tale}}
\def\etale{{\acute{e}tale}}
\def\QCoh{\mathit{QCoh}}
\def\Ker{\mathop{\mathrm{Ker}}}
\def\Im{\mathop{\mathrm{Im}}}
\def\Coker{\mathop{\mathrm{Coker}}}
\def\Coim{\mathop{\mathrm{Coim}}}

% Boxtimes
%
\DeclareMathSymbol{\boxtimes}{\mathbin}{AMSa}{"02}

%
% Macros for moduli stacks/spaces
%
\def\QCohstack{\mathcal{QC}\!\mathit{oh}}
\def\Cohstack{\mathcal{C}\!\mathit{oh}}
\def\Spacesstack{\mathcal{S}\!\mathit{paces}}
\def\Quotfunctor{\mathrm{Quot}}
\def\Hilbfunctor{\mathrm{Hilb}}
\def\Curvesstack{\mathcal{C}\!\mathit{urves}}
\def\Polarizedstack{\mathcal{P}\!\mathit{olarized}}
\def\Complexesstack{\mathcal{C}\!\mathit{omplexes}}
% \Pic is the operator that assigns to X its picard group, usage \Pic(X)
% \Picardstack_{X/B} denotes the Picard stack of X over B
% \Picardfunctor_{X/B} denotes the Picard functor of X over B
\def\Pic{\mathop{\mathrm{Pic}}\nolimits}
\def\Picardstack{\mathcal{P}\!\mathit{ic}}
\def\Picardfunctor{\mathrm{Pic}}
\def\Deformationcategory{\mathcal{D}\!\mathit{ef}}
\usepackage{cleveref}
\crefname{originalreference}{original reference}{original references}
\Crefname{originalreference}{Original reference}{Original references}
\makeatletter
\newlabel{lemma-m-adic-topology}{{06SE}{source}{}{sourcepointer-1}{}}
\newlabel{lemma-m-adic-topology@cref}{{[originalreference][0][]06SE}{[1][0][]source}{}{}{}}
\makeatother
\crefname{originalreference}{original reference}{original references}
\Crefname{originalreference}{Original reference}{Original references}
\makeatletter
\newlabel{lemma-surjective-cotangent-space}{{06GZ}{source}{}{sourcepointer-2}{}}
\newlabel{lemma-surjective-cotangent-space@cref}{{[originalreference][0][]06GZ}{[1][0][]source}{}{}{}}
\makeatother
\title{A bounded inverse system of Artinian algebras has a limit in the completed category, with the kernels inducing the adic topology}
\author{The Stacks Project Authors}
\date{stacks-project:formal-defos.tex:local-snapshot-20261012; GNU FDL 1.2 or later}
\begin{document}
\maketitle
\noindent\textit{Editorial note.} All delivered mathematical text is the author's own, copied from the pinned Stacks Project chapter file formal-defos.tex; no statement or proof step was rewritten, reordered or removed.\par\medskip
\noindent\textit{Editorial note.} The counted proof is the author's complete argument: the stabilization of the relative cotangent space dimensions and the resulting compatible polynomial presentation, the passage to the adic completion with the map onto the system of quotients, and the application of the topology lemma identifying the limit and the ideals defining its topology.\par\medskip
\noindent\textit{Editorial note.} The label lemma-limit-artinian names the counted statement, and the proof environment immediately following it in the source is delivered verbatim as its proof.\par\medskip
\noindent\textit{Editorial note.} The counted statement is carried by the author's own lemma environment in the chapter source.\par\medskip
\noindent\textit{Editorial note.} Pre-open gate: python operations/agent-authoring/tools/precard\_gate.py --source formal-defos.tex --label lemma-limit-artinian --start 1277 --end 1313 --local-id 28020 returned RESULT=OK (L28020 free, occupancy FREE, public interval clear). Candidate taken from sources/stacks-formal-defos/candidates-halfA.json.\par\medskip
\noindent\textit{Editorial note.} The statement and proof are copied from the pinned chapter file; the cited inputs are kept as external pointers with their canonical Stacks tags.\par\medskip
\noindent\textit{Allowed standard prerequisites (editorial).}\begin{itemize}
\item Surjectivity on relative cotangent spaces. Same chapter, the lemma that a surjective map of Artinian algebras induces a surjection on the relative cotangent spaces. Supplies the surjectivity of the transition maps on the cotangent spaces, hence the stabilization of their dimension.
\item The adic topology of a quotient of a complete Noetherian algebra. Same chapter, the lemma describing the quotient of a complete local algebra by the intersection of a compatible system of ideals as the limit of the quotients, with the induced topology. Supplies the identification R/J = lim A\_n and the statement about the topology.
\end{itemize}\medskip
\ifdefined\hypertarget\hypertarget{sourcepointer-1}{}\fi
\noindent\textit{Source reference.} Original reference 06SE: sources/stacks-formal-defos/pinned-original-source.tex, source lines 1251--1251. Prerequisite pointer invoked at proof line(s) 1307; it names a standard fact of the same chapter that the counted argument uses, and every step of the counted proof built on it is retained verbatim in the delivered body.\par\medskip
\ifdefined\hypertarget\hypertarget{sourcepointer-2}{}\fi
\noindent\textit{Source reference.} Original reference 06GZ: sources/stacks-formal-defos/pinned-original-source.tex, source lines 1086--1086. Prerequisite pointer invoked at proof line(s) 1292; it names a standard fact of the same chapter that the counted argument uses, and every step of the counted proof built on it is retained verbatim in the delivered body.\par\medskip
\begin{lemma}
\label{lemma-limit-artinian}
Let $\ldots \to A_3 \to A_2 \to A_1$ be a sequence of surjective
ring maps in $\mathcal{C}_\Lambda$. If
$\dim_k (\mathfrak m_{A_n}/\mathfrak m_{A_n}^2)$ is bounded, then
$S = \lim A_n$ is an object in $\widehat{\mathcal{C}}_\Lambda$
and the ideals $I_n = \Ker(S \to A_n)$ define the
$\mathfrak m_S$-adic topology on $S$.
\end{lemma}

\begin{proof}
We will use freely that the maps $S \to A_n$ are surjective for all $n$.
Note that the maps
$\mathfrak m_{A_{n + 1}}/\mathfrak m_{A_{n + 1}}^2 \to
\mathfrak m_{A_n}/\mathfrak m_{A_n}^2$ are surjective, see
Lemma \ref{lemma-surjective-cotangent-space}.
Hence for $n$ sufficiently large the dimension
$\dim_k (\mathfrak m_{A_n}/\mathfrak m_{A_n}^2)$ stabilizes to an
integer, say $r$.
Thus we can find $x_1, \ldots, x_r \in \mathfrak m_S$ whose images in
$A_n$ generate $\mathfrak m_{A_n}$. Moreover, pick $y_1, \ldots, y_t \in S$
whose images in $k$ generate $k$ over $\Lambda$. Then we get a ring map
$P = \Lambda[z_1, \ldots, z_{r + t}] \to S$, $z_i \mapsto x_i$ and
$z_{r + j} \mapsto y_j$ such that the composition
$P \to S \to A_n$ is surjective for all $n$. Let $\mathfrak m \subset P$
be the kernel of $P \to k$. Let $R = P^\wedge$ be the $\mathfrak m$-adic
completion of $P$; this is an object of $\widehat{\mathcal{C}}_\Lambda$.
Since we still have the compatible system of (surjective) maps $R \to A_n$
we get a map $R \to S$. Set $J_n = \Ker(R \to A_n)$.
Set $J = \bigcap J_n$. By
Lemma \ref{lemma-m-adic-topology}
we see that $R/J = \lim R/J_n = \lim A_n = S$
and that the ideals $J_n/J = I_n$ define the $\mathfrak m$-adic topology.
(Note that for each $n$ we have $\mathfrak m_R^{N_n} \subset J_n$ for
some $N_n$ and not necessarily $N_n = n$, so a renumbering of the ideals
$J_n$ may be necessary before applying the lemma.)
\end{proof}

\end{document}
